% !PRINTED 433
\weberpage{433}
\weq{1}{\cos n\varphi + i\sin n\varphi = (\cos\varphi + i\sin\varphi)^n,}
and if we apply the binomial theorem on the right-hand side, then by separation of the real from the imaginary it follows that
\[
\begin{aligned}
\cos n\varphi &= \cos^n\varphi - B_2^{(n)}\cos^{n-2}\varphi\sin^2\varphi\\
&\quad + B_4^{(n)}\cos^{n-4}\varphi\sin^4\varphi - \cdots
\end{aligned}
\]
\weq{2}{
\begin{aligned}
\frac{\sin n\varphi}{\sin\varphi} &= n\cos^{n-1}\varphi - B_3^{(n)}\cos^{n-3}\varphi\sin^2\varphi\\
&\quad + B_5^{(n)}\cos^{n-5}\varphi\sin^4\varphi - \cdots
\end{aligned}
}
wherein $B_\nu^{(n)}$ are, as before, the binomial coefficients, and the sums on the right-hand side are to be continued as far as no negative exponents of $\cos\varphi$ occur. Since now $\sin^2\varphi = 1 - \cos^2\varphi$, these two expressions can be represented rationally in terms of $\cos\varphi$; we set
\weq{3}{2\cos\varphi = x,}
and introduce the notation
\weq{4}{2\cos n\varphi = A_n(x), \qquad \frac{\sin n\varphi}{\sin\varphi} = B_n(x),}
wherein, then, $A_n(x)$ and $B_n(x)$ are integral rational functions of $x$, $A_n(x)$ of degree $n$, $B_n(x)$ of degree $n-1$. We wish to determine the general law of these functions. We set them up for the first values $n = 1, 2, 3\ldots$, and then derive a recursion formula for the computation of the higher functions. One finds
\weq{5}{
\begin{aligned}
A_1(x) &= x & B_1(x) &= 1\\
A_2(x) &= x^2 - 2, & B_2(x) &= x\\
A_3(x) &= x^3 - 3x, & B_3(x) &= x^2 - 1.
\end{aligned}
}
But now, by well-known formulae,
\[
\begin{aligned}
\cos(n+1)\varphi + \cos(n-1)\varphi &= 2\cos\varphi\cos n\varphi\\
\sin(n+1)\varphi + \sin(n-1)\varphi &= 2\cos\varphi\sin n\varphi,
\end{aligned}
\]
from which the recursion formulae for $A_n$ and $B_n$ follow:
\weq{6}{
\begin{aligned}
A_{n+1}(x) &= x A_n(x) - A_{n-1}(x),\\
B_{n+1}(x) &= x B_n(x) - B_{n-1}(x).
\end{aligned}
}
From this one can compute $A_n$, $B_n$ for any arbitrary $n$, if these functions are known for $n = 1$ and $n = 2$. Thus one finds
%CONT>>
